\begin{document}

\title{$\neexp \subseteq \mip^*$\vspace*{10pt}}

\author{Anand Natarajan\thanks{anandn@caltech.edu}\\
 \small{\sl California Institute of Technology} \and John Wright\thanks{jswright@mit.edu}\\ \small{\sl Massachusetts Institute of Technology}\vspace*{10pt}
}

\date{\today}

\maketitle


%%%%%%%%%%% jump to abstract
%%%---------- open: abstract
\begin{abstract}
We study multiprover interactive proof systems.
The power of classical multiprover interactive proof systems,
in which the provers do not share entanglement,
was characterized in a famous work by Babai, Fortnow, and Lund
(Computational Complexity 1991),
whose main result was the equality $\MIP = \NEXP$.
The power of quantum multiprover interactive proof systems,
in which the provers are allowed to share entanglement,
has proven to be much more difficult to characterize.
The best known lower-bound on $\mip^*$ is $\nexp \subseteq \mip^*$ due
to Ito and Vidick~(FOCS 2012).
As for upper bounds, $\mip^*$ could be as large as $\mathsf{RE}$, the class of recursively enumerable languages.

The main result of this work is the inclusion $\neexp = \ntime[2^{2^{\poly(n)}}] \subseteq \mip^*$.
This is an exponential improvement over the prior lower bound
and shows that proof systems with entangled provers are at least exponentially more
powerful than classical provers. In our protocol the verifier
delegates a classical, exponentially large $\MIP$ protocol for
$\neexp$ to two entangled provers: the provers obtain their
exponentially large questions by measuring their shared state, and use
a classical PCP to certify the correctness of their exponentially-long
answers. For the soundness of our protocol, it is
crucial that each player should not only sample its own question correctly
but also avoid performing measurements that would reveal the
\emph{other} player's sampled question. We ensure this by commanding
the players to perform a complementary measurement, relying on the Heisenberg uncertainty
principle to prevent the forbidden measurements from being performed.



% We also ``scale down" this result to show that $\NEXP \subseteq \mip^*[\log(n), 1]$,
% improving on the recent  games quantum PCP theorem by Natarajan and Vidick~\cite{NV18b}.
\end{abstract}

%%% Local Variables:
%%% mode: latex
%%% TeX-master: "main"
%%% End:
%%%---------- close: abstract

\newpage

\tableofcontents
\vfill
\thispagestyle{empty}
\newpage


\part{Introduction}

%%%%%%%%%%% jump to introduction
%%%---------- open: introduction
%!TEX root = main.tex

